\documentclass[11pt]{letter}
\usepackage[margin=2.5cm]{geometry}
\usepackage{mathpazo}
\usepackage{url}
\usepackage{hyperref}
\hypersetup{colorlinks=true,urlcolor=blue,linkcolor=blue}

\date{April 27, 2026}

\signature{Yan Ma$^{1,3}$ \\ Lizhuo Zhang$^{2}$ \\ $^{1}$Hunan Agricultural University, Changsha 410128, China \\ $^{2}$School of Information and Intelligence, Hunan Agricultural University \\ $^{3}$Changsha University of Science and Technology, Changsha 410114, China \\ Corresponding author e-mail: zhanglizhuo@hunau.edu.cn}

\address{Hunan Agricultural University \\ Changsha 410128, China}

\begin{document}

\begin{letter}{Editorial Office \\ \textit{Electronics} \\ MDPI}

\opening{Dear Editor,}

We would like to submit the enclosed manuscript, ``Benchmarking CLIP-Family Models for Zero-Shot Classroom Activity Recognition: Prompt and Metric Choice Matter,'' for consideration as an original research article in \textit{Electronics}.

\textbf{Background and motivation.}
Zero-shot vision--language models are increasingly used as practical baselines for domain transfer, yet their behavior on classroom activity recognition remains underexplored. Existing evaluations typically emphasize headline top-1 scores without testing whether prompt choice and metric choice materially change the conclusions drawn from benchmark results, especially under class imbalance and multi-label structure.

\textbf{Summary of contributions.}
This study makes four main contributions. First, it constructs a multi-scene classroom benchmark from three public SCB subsets covering teacher behavior, student learning actions, and subtle micro-actions. Second, it provides a unified probe-based evaluation of five CLIP-family backbones under four uniform prompt strategies and a training-free class-aware prompt ensemble (CAPE), showing that prompt design is often more consequential than backbone choice and that CAPE is task-dependent rather than universally optimal; a three-principle ablation further attributes CAPE's gains to visual grounding and class-discriminative contrast rather than to averaging more prompts. Third, it demonstrates that evaluation protocol matters: a supervised linear-probe baseline reaches Sample-F1 of 88--90\%, Macro-F1 of 68--78\%, and Micro-F1 of 89--91\% on the multi-label TeacherBehavior subset, while zero-shot CAPE only reaches 60--66\%, 49--61\%, and 61--70\%, exposing the headline Hit@1 advantage as a lenient-metric artifact; the same lenient-metric caveat applies to BowTurnHead, where the multi-label majority floor is already 90.69\% so the nominal 93.27\% headline reflects a $\sim$2.6~pp effect on top of a near-saturated metric, and the stricter single-label primary-accuracy view (68.3\% [64.6, 72.1]) is the more reliable evidence of CAPE's benefit on micro-actions. Fourth, additional cross-family checks against three Ollama-deployed MLLMs and a direct comparison between EVA02 under fixed CAPE and EVA02 under its best prompt strategy show that benchmark conclusions in classroom zero-shot recognition should rely on multiple metrics and explicitly disclosed prompt strategies rather than isolated headline numbers.

\textbf{Relevance to \textit{Electronics}.}
The manuscript fits the journal's interests in intelligent sensing, computer vision, and AI-enabled educational systems. Beyond reporting benchmark scores, it addresses a practical methodological issue relevant to deployable vision systems: how prompt design and metric choice affect the reliability of zero-shot conclusions in complex real-world scenes.

\textbf{Statements.}
This manuscript has not been published previously and is not under consideration by any other journal. All authors have read and approved the manuscript. The study reuses only existing public dataset releases and does not involve new human-subject data collection. The authors declare no conflicts of interest. This work was funded by the Education Department of Hunan Province, China (Project Number: HNJG-2022-0608).

The SCB datasets analyzed in this study are publicly available from their original maintainers at \url{https://huggingface.co/datasets/wintonYF/SCB-Dataset}. The reproducibility package, including experiment code, prompt templates, and released summary result files, is available at \url{https://github.com/zhanglizhuo/scb-clip-benchmark}.

We believe the manuscript will be of interest to readers working on vision--language models, zero-shot learning, classroom analytics, and trustworthy evaluation of domain-specific AI systems. We appreciate your consideration.

\closing{Sincerely,}

\end{letter}
\end{document}